
\Needspace{12\baselineskip}
\subsubsection*{Naučnoistraživački rad -- \ClanCetiri}

\textbf{Naslov rada:}
\emph{A machine learning based credit card fraud detection using
the GA algorithm for feature selection}.\\

\textbf{Autori:}
Emmanuel Ileberi, Yanxia Sun i Zenghui Wang.\\

\textbf{Izvor:}
\emph{Journal of Big Data}, vol. 9, članak 24,
2022. \cite{ileberi2022fraud}

\begin{table}[H]
\centering
\caption{Sažeti pregled naučnoistraživačkog rada Ileberija i saradnika (2022)}
\label{tab:ileberi-pregled}
\small
\renewcommand{\arraystretch}{1.22}
\begin{tabularx}{\linewidth}{@{}p{0.25\linewidth}Y@{}}
\toprule
\textbf{Stavka} & \textbf{Opis} \\
\midrule

Problem &
Detekcija prevarnih transakcija kreditnim karticama uz izbor
optimalnog podskupa atributa pomoću genetskog algoritma (GA). \\

Skup podataka &
Evropski skup transakcija kreditnim karticama (Kaggle, septembar 2013.):
284\,807 transakcija, 0,172\,\% prevara, 30 numeričkih atributa
(V1--V28 dobijeni PCA transformacijom, \texttt{Time}, \texttt{Amount}).
Dodatna provjera na sintetičkom skupu s oko 24,4 miliona transakcija. \\

Korištene metode &
Genetski algoritam za selekciju atributa (funkcija prilagođenosti je
tačnost Random Foresta); klasifikatori: stablo odlučivanja, Random Forest,
logistička regresija, ANN i Naive Bayes; SMOTE za uravnotežavanje
i \emph{min-max} normalizacija. \\

Metrike &
Tačnost (glavna metrika), osjetljivost (recall), preciznost,
F1-mjera i AUC. \\

Najbolji rezultat &
\textbf{GA-RF} s vektorom atributa v5 (13 atributa): tačnost
\textbf{99,98\,\%}, preciznost 95,34\,\%, recall 72,56\,\%,
F1-mjera 82,41\,\%. \\

\bottomrule
\end{tabularx}
\end{table}

\textbf{Cilj i istraživački problem.}
Autori smatraju da veliki broj atributa može narušiti performanse
klasifikatora i predlažu sistem u kojem genetski algoritam bira
podskup atributa, a zatim se na tom podskupu trenira i upoređuje
pet klasičnih algoritama mašinskog učenja.

\textbf{Skup podataka i priprema.}
Korišten je najpoznatiji javni skup za ovaj problem: transakcije
evropskih korisnika kartica tokom dva dana, sa samo 492 prevare.
Za razliku od našeg skupa, svi atributi su numerički i anonimizovani.
Klasna nebalansiranost je ublažena metodom SMOTE, a atributi su
normalizovani \emph{min-max} postupkom. Genetski algoritam je
pronašao pet kandidatskih vektora atributa (v1--v5) dužine od 9 do
18 atributa, a atribut \texttt{Amount} se nalazi u svakom od njih.

\textbf{Poređenje eksperimentalnih rezultata.}
Tabela~\ref{tab:ileberi-rezultati} prikazuje rezultate za vektor v5
(Tabela 6 rada) i za vektor v1, na kojem je ostvarena najveća F1-mjera
(Tabela 2 rada).

\begin{table}[H]
\centering
\caption{Rezultati klasifikatora iz rada Ileberija i saradnika (u \%)}
\label{tab:ileberi-rezultati}
\small
\setlength{\tabcolsep}{4pt}
\renewcommand{\arraystretch}{1.18}
\begin{tabularx}{\linewidth}{@{}Y c c c c@{}}
\toprule
\textbf{Model (vektor atributa)} & \textbf{Acc.} & \textbf{Recall} & \textbf{Prec.} & \textbf{F1} \\
\midrule
\textbf{Random Forest (v5)} & \textbf{99,98} & \textbf{72,56} & \textbf{95,34} & \textbf{82,41} \\
Stablo odlučivanja (v5) & 99,89 & 72,56 & 65,07 & 68,61 \\
ANN (v5) & 99,08 & 77,87 & 12,27 & 21,20 \\
Naive Bayes (v5) & 99,44 & 57,52 & 15,85 & 24,85 \\
Logistička regresija (v5) & 99,77 & 46,90 & 34,64 & 39,84 \\
\midrule
Random Forest (v1) & 99,94 & 76,99 & 89,69 & 82,85 \\
ANN (v1) & 99,94 & 77,87 & 84,61 & 81,10 \\
Naive Bayes (v1) & 98,13 & 84,95 & 6,83 & 12,65 \\
\bottomrule
\end{tabularx}
\end{table}

\textbf{Tumačenje rezultata.}
Random Forest je najbolji ili među najboljima za svaki vektor atributa.
Iako svi modeli imaju tačnost iznad 98\,\%, preciznost i F1-mjera se
drastično razlikuju: Naive Bayes s vektorom v1 ima tačnost 98,13\,\%,
a preciznost samo 6,83\,\%, što znači da je velika većina transakcija
koje označi kao prevarne zapravo legitimna. Rezultati pokazuju i da
izbor atributa snažno utiče na pojedine modele (npr. F1 ANN-a pada
s 81,10\,\% na 21,20\,\% između vektora v1 i v5). Treba imati na umu
da autori kao glavnu metriku koriste tačnost, koja je na skupu s
0,172\,\% prevara slabo informativna.

\textbf{Zaključak i značaj za istraživanje.}
Rad potvrđuje da je Random Forest pouzdan izbor za detekciju prevara
i da Naive Bayes ima najlošiji odnos preciznosti i osjetljivosti.
Za naš projekat je važna i pouka o evaluaciji: kod nebalansiranih
podataka rezultate treba prikazivati kroz preciznost, recall i
F1-mjeru za klasu prevare, a ne samo kroz tačnost.
